\section{Closing the signed-increment bootstrap}\label{sec:closure}
We now prove Proposition~\ref{prop:increment}. The preceding estimates
were conditional on the signed bound \eqref{eq:bootstrap}. Their
constants permit a strict improvement, uniformly over the chosen
compact solution interval. A continuity argument then removes that
assumption.

Fix an interval $[t_1,t_2]$ of length $I$ on which
$w/8\le q_X\le8w$. Write the full bootstrap bound as
\[
 \mathcal B_w(I)=MP\sqrt I+A I S(w).
\]
If $\mathcal B_w(I)\ge32w$, then
\[
 |V_X(t_2)-V_X(t_1)|\le q_X(t_2)+q_X(t_1)
 \le16w\le\tfrac12\mathcal B_w(I).
\]
Otherwise, once $M\ge32$, we have $I\le(w/P)^2$, as required in
the estimates. The case $I=0$ is immediate.

Let $0<\eta<1/2$ and choose a Lipschitz function $\psi$ on
$[t_1,t_2]$ that vanishes at the endpoints, equals one off the two
edge intervals of length $\eta I$, and has linear monotone ramps on those
intervals. The weighted impulse estimate \eqref{sel:bulk} gives
\begin{equation}\label{eq:closure-bulk}
 \left|\int_{t_1}^{t_2}\psi(t)K(t)\,dt\right|
 \le C_\eta\bigl[P\sqrt I+(M+\sqrt A)I S(w)\bigr].
\end{equation}
We restore each edge using the signed bootstrap, without estimating
the absolute force there. On the left edge put $\chi=1-\psi$ and
$F(t)=V_X(t)-V_X(t_1)$. Since $F(t_1)=0$ and $\chi$ vanishes
at the other edge endpoint,
\[
 \int\chi(t)K(t)\,dt=-\int\chi'(t)F(t)\,dt.
\]
The monotonicity of the ramp gives $\int|\chi'|=1$; every such
$F(t)$ has norm at most $\mathcal B_w(\eta I)$ by
\eqref{eq:bootstrap}. The right edge has the same bound, anchoring
at $t_2$. Hence
\begin{equation}\label{eq:closure-improvement}
 |V_X(t_2)-V_X(t_1)|
 \le (C_\eta+2M\sqrt\eta)P\sqrt I
  +\bigl[C_\eta(M+\sqrt A)+2A\eta\bigr]I S(w).
\end{equation}

Choose $\eta$ so that $2\sqrt\eta\le1/8$ and $2\eta\le1/8$.
Fix $M\ge\max(32,8C_\eta)$ and then
\[
 A\ge\max(1,16C_\eta M,256C_\eta^2).
\]
The right side of \eqref{eq:closure-improvement} is at most
$\mathcal B_w(I)/4$, and in particular at most
$\mathcal B_w(I)/2$. Finally choose $P_0$ large enough for all
preceding estimates with these fixed $M,A$. These choices depend
only on the datum and $T_0$, and not on $t_*$.
Thus every bootstrap interval satisfies the uniform improvement
\begin{equation}\label{eq:strict-improvement}
 |V_X(t_2)-V_X(t_1)|\le\tfrac12\mathcal B_w(t_2-t_1).
\end{equation}

For completeness, fix $P\ge P_0$ satisfying \eqref{eq:P-range} on
$[0,t_*]$. Let $\mathcal T$ be the set of $\tau\in[0,t_*]$ for
which \eqref{eq:bootstrap} holds for all eligible receivers, scales
and intervals contained in $[0,\tau]$. It contains zero and is an
initial interval. It is closed: truncate any nontrivial eligible
interval at an increasing sequence of endpoints from below and use
continuity of $V_X$. The range condition is preserved by truncation.

It is also relatively open to the right. On this fixed compact
solution interval there is a finite bound $C_*$ for $|K|$ on all
supported characteristics. Consequently some $\ell_0>0$ ensures
\eqref{eq:bootstrap} for every interval of length at most $\ell_0$,
since $C_*I\le MP\sqrt I$ there. If $\tau\in\mathcal T$, a longer
eligible interval ending in $[\tau,\tau+\varepsilon]$ has a prefix
ending at $\tau$. The strict estimate \eqref{eq:strict-improvement}
on that prefix is at most half the full interval's bound. The suffix
has impulse at most $C_*\varepsilon$, which is smaller than the
remaining margin when $\varepsilon$ is sufficiently small, uniformly
because the full bound is at least $MP\sqrt{\ell_0}$. An interval
starting after $\tau$ is short if $\varepsilon<\ell_0$.
This proves relative openness. Hence $\mathcal T=[0,t_*]$ and
Proposition~\ref{prop:increment} follows.

The compact-time force bound $C_*$ was used only to propagate the
bootstrap for a fixed $P$. It does not enter $M,A$ or $P_0$.
This distinction preserves the uniformity needed in the
momentum-doubling argument of Section~\ref{sec:setup}.
